\documentclass{article}
\begin{document}
\abstract[Abstract]{We bound the difference. The summary measures stayed at chance wherever discordance lay in distributional shape, and the distance recovered the structure in every world we built, with coverage near nominal across all seven scenarios of the study design that we report in the simulation section below.}
\section{Introduction}\label{sec:intro}
\begin{table}\caption{defined before cited\label{tab:early}}\end{table}
Table~\ref{tab:early} is defined above its first citation. The distance has a unique theoretical connection to treatment effect heterogeneity that other measures lack entirely.
\begin{figure}\caption{never cited\label{fig:orphan}}\end{figure}
\section{Methods}\label{sec:methods}
The bound is unique to the distance and specific to it alone. Section 1 states that the distance has a unique theoretical connection to treatment effect heterogeneity that other measures lack entirely.
\section{Results}\label{sec:results}
Coverage was 0.70 at $n=50$, 0.92 at $n=100$ and 0.95 at $n=200$ across 10,000 replications, as Supplement A and Supplement C show.
\section{Discussion}\label{sec:discussion}
Anchor-wise screening violated the pool-wide requirement in 61 to 68 percent of replications across 100,000 runs, the residual breach fell from 16 to below 1 percent between 25 and 400 patients, and the 3 cells of Set 2 agreed with the 4 cells of Set 1 at 0.636 and 0.533 and 0.447, which is a simulation reported in the Discussion rather than in the Results section where evidence belongs by convention. Section 1 states that the distance has a unique theoretical connection to treatment effect heterogeneity that other measures lack entirely.
\end{document}
